\section{Как кластерная модель меняет положение бинодали}
\label{sec:ce-phases}

\subsection{Расчёт фазовых свободных энергий}
Для каждого набора эффективных кластерных взаимодействий (ECI) были
построены три термодинамические модели. В статической модели при опорном
объёме \(v_0=17{,}28833\)~\AA$^3$/атом химические заполнения имеют
больцмановские веса, определяемые релаксированными статическими энергиями.
В гармоническом приближении (HA) к ним добавляется квантовая колебательная
свободная энергия при том же \(v_0\). В квазигармоническом приближении
(QHA) статические и колебательные взаимодействия зависят от удельного
объёма \(v\), и объём выбирается в равновесии с внешним давлением \(p=0\).
Отношение осей остаётся равным \(c/a=1{,}56374\) во всех трёх моделях.
В QHA конфигурация уже несёт свободную энергию колебаний для данного
\(T,v\), так что в решёточном Монте-Карло изменяются заполнения и
общий объём.

В приближении среднего поля (MFA) переменная \(q\) задаёт долю Al
на предпочтительной подрешётке, а \(r=(4x-q)/3\) --- долю на каждой
из остальных трёх. При \(q=r=x\) раствор разупорядочен. Для заданного \(x\)
произведения спинов усредняли как произведения одноузельных средних,
а затем минимизировали свободную энергию на атом
\[
 f_{\rm MFA}(x,q,v,T)=e_{\rm CE}^{\rm MFA}(q,r,v,T)
       +\frac{k_{\rm B}T}{4}\big[h(q)+3h(r)\big],
 \qquad h(y)=y\ln y+(1-y)\ln(1-y).
\]
Минимум искали по \(x\le q\le\min(4x,1)\) и, для QHA, по \(v\).
Общая касательная к полученной зависимости от \(x\) даёт обе границы.
В этой схеме дальний порядок может сформироваться за счёт различия
подрешёток, а корреляции между заполнениями соседних узлов отброшены.

В решёточном Монте-Карло корреляции возникают в ходе замены Ti на Al
и обратно. Чтобы измерить изотерму между двумя устойчивыми ветвями,
в окне \(l\) к энергии ячейки помимо поля \(-N\Delta\mu_{0,l}x\)
добавляли \(Nk_{\rm B}T\kappa_l(x-x_{0,l})^2\). Это полубольшой
канонический ансамбль с ограниченной дисперсией состава (VCSGC)~\cite{sadigh2012}.
По фактическому среднему составу каждого окна определяли производную
фазовой свободной энергии
\[
 \widehat{\Delta\mu}_l=\Delta\mu_{0,l}
              -2k_{\rm B}T\kappa_l\big(\langle x\rangle_l-x_{0,l}\big).
\]
Пары \((\langle x\rangle_l,\widehat{\Delta\mu}_l)\)
интерполировали сохраняющим форму кусочно-кубическим полиномом PCHIP,
интегрировали по \(x\) и искали общую касательную. Для одного интервала
эта операция эквивалентна условию равных площадей
\(\int_{x_\alpha}^{x_{\alpha_2}}[\widehat{\Delta\mu}(x)-\Delta\mu^*]\,dx=0\)
при \(\widehat{\Delta\mu}(x_\alpha)=\widehat{\Delta\mu}(x_{\alpha_2})=\Delta\mu^*\).
Интегрируется производная по составу уже заданного решёточного гамильтониана.
В разделе~\ref{sec:smallcells} приведён другой интеграл, по параметру
переключения между двумя атомистическими гамильтонианами.

\subsection{Сопоставление шести расчётов при 1000 К}
В табл.~\ref{tab:ce-six} сопоставлены авторские коэффициенты C46 и
коэффициенты, заново полученные на независимо отобранных структурах ATAT.
Для обеих QHA--MFA использован одинаковый кубический сплайн по объёму
с граничным условием \textit{not-a-knot}. Для MC один и тот же алгоритм
уточнения касательной применён к исходным наборам средних по окнам.
В авторских сечениях на 5120 узлах использованы соответственно 15, 12 и
13 окон для статики, HA и QHA, в независимых --- по 35 окон.

\begin{table}[htbp]
\centering\small
\caption{Составы фаз при 1000 К, ат.\% Al. Последние два столбца дают
изменение при переходе от авторских ECI к независимой подгонке, п.п.}
\label{tab:ce-six}
\begin{tabular}{@{}llrr rr rr@{}}
\toprule
 & & \multicolumn{2}{c}{Авторский C46} &
 \multicolumn{2}{c}{Набор ATAT} & \multicolumn{2}{c}{ATAT минус C46}\\
\cmidrule(lr){3-4}\cmidrule(lr){5-6}\cmidrule(l){7-8}
Модель & Метод & \(\alpha\) & \(\alpha_2\) & \(\alpha\) & \(\alpha_2\)
 & \(\alpha\) & \(\alpha_2\)\\
\midrule
Статика & MFA & 6,159 & 20,799 & 6,758 & 22,593 & \(+0{,}599\) & \(+1{,}795\)\\
HA      & MFA & 8,926 & 19,421 & 9,653 & 21,079 & \(+0{,}727\) & \(+1{,}658\)\\
QHA     & MFA & 8,789 & 21,308 & 9,431 & 22,413 & \(+0{,}643\) & \(+1{,}106\)\\
\addlinespace
Статика & MC  & 11,447 & 20,472 & 11,120 & 22,290 & \(-0{,}328\) & \(+1{,}818\)\\
HA      & MC  & 15,281 & 20,612 & 15,173 & 21,656 & \(-0{,}109\) & \(+1{,}043\)\\
QHA     & MC  & 14,958 & 21,958 & 14,764 & 22,707 & \(-0{,}194\) & \(+0{,}748\)\\
\bottomrule
\end{tabular}
\end{table}

\paragraph{Химические корреляции.}
При неизменных авторских QHA-коэффициентах переход от MFA к MC смещает
границы на \(+6{,}169\) и \(+0{,}651\) п.п. Ширина двухфазной области
уменьшается с \(12{,}519\) до \(7{,}000\) п.п. Особенно чувствителен
неупорядоченный раствор. Независимые заполнения в MFA лишают его
возможности снизить свободную энергию за счёт ближнего химического
порядка. Направление наблюдаемого сдвига согласуется со стабилизацией
раствора, хотя сам вклад в свободную энергию нельзя получить вычитанием
двух составов. Независимая подгонка также даёт большой сдвиг границы
\(\alpha\), \(+5{,}333\) п.п.

\paragraph{Колебания и объём.}
В авторском MC добавление гармонической свободной энергии при фиксированном
объёме увеличивает \(x_\alpha\) на \(3{,}834\) п.п., а \(x_{\alpha_2}\)
на \(0{,}140\) п.п. Ширина уменьшается с \(9{,}025\) до \(5{,}331\) п.п.
Переход HA\(\to\)QHA затем даёт \(-0{,}323\) и \(+1{,}346\) п.п.
и увеличивает ширину до \(7{,}000\) п.п. Таким образом, при фиксированном
объёме сильнее меняется растворимость Al в \(\alpha\), а при разрешении
объёмной релаксации --- состав упорядоченной фазы. В авторском QHA--MFA
на собственных контактах объёмы составили \(17{,}606\) и
\(17{,}267\)~\AA$^3$/атом. Замена граничного условия объёмного сплайна
на естественное меняет правый MFA-контакт всего на \(-0{,}024\) п.п.

\paragraph{Почему расходятся две подгонки.}
При одинаковом симметрийном базисе независимая подгонка во всех шести
сравнениях переносит правую границу к большему содержанию Al.
В QHA--MC изменение \(+0{,}748\) п.п. существенно превышает расхождение
двух реализаций поиска касательной. На идентичных окнах для всех
разрешённых авторских сечений максимальная разница обработчиков
составила \(0{,}00244\) п.п. На независимом наборе она также
меньше \(0{,}0025\) п.п. для сопоставленных наиболее широких интервалов.

В независимой QHA при 1000 К серия из 18 окон на 768 узлах дала
\(14{,}431/22{,}770\) ат.\% Al. Серия из 35 окон на 5120 узлах
дала \(14{,}764/22{,}707\) ат.\%. Здесь одновременно менялись размер
и сетка окон, поэтому сдвиг \(+0{,}334/-0{,}063\) п.п. характеризует
их совместное действие. В большом сечении наблюдаемые гистограммы
состава образовали 35 несвязанных компонент. Числа в таблице получены
интегрированием средних химических потенциалов. Из этих гистограмм
восстановить относительные веса всех окон методом WHAM не удалось.

Чувствительность к отдельным взаимодействиям измерена для MFA.
Вдоль прямого пути от авторских к собственным \(J(V)\) интегрировали
производную правого QHA-контакта по каждому коэффициенту.
В изменение \(+1{,}106\) п.п. треугольники \(t_{334}\) и \(t_{123}\)
вносят \(+0{,}918\) и \(+0{,}709\) п.п., а одноузельный и постоянный
члены дают \(-0{,}410\) и \(+0{,}122\) п.п. При фиксированном \(v\)
два последних члена аффинны по \(x\) и сохраняют составы касания.
При объёмной релаксации их зависимость от \(v\) меняет равновесный объём
и относительную свободную энергию ветвей. Конечная замена только
\(t_{334}(V)\) и \(t_{123}(V)\) в C46 переносит правый контакт
с \(21{,}308\) до \(22{,}813\) ат.\%. Остальные изменения ECI частично
компенсируют этот сдвиг. Разбиение по орбитам зависит от пути замены
коэффициентов и относится к свободной энергии среднего поля.

Практическое следствие видно при сравнении с прямым MTP при том же
\(c/a\) (раздел~\ref{sec:bulk-c46}). Прямые составы равны
\(16{,}083\) и \(22{,}159\) ат.\%. Разность «прямой MTP минус
CE--QHA--MC» для C46 составляет \(+1{,}125\) и \(+0{,}201\) п.п.,
а для независимой подгонки \(+1{,}318\) и \(-0{,}547\) п.п.
Обе подгонки недооценивают растворимость Al в \(\alpha\). Для
\(\alpha_2\) смена обучающей выборки меняет даже знак отклонения
от прямого классического расчёта. Здесь, помимо точности ECI, остаётся
различие в колебательной модели и статистике ядер.

\subsection{Температурная зависимость и разрешение узкого интервала}
% FIG_CE_TEMPERATURE
В авторском QHA--MC при охлаждении до 600 К область сосуществования
расширяется до \(4{,}59\)--\(24{,}09\) ат.\% Al (рис.~\ref{fig:ce-temp}).
При 1000 К её ширина составляет \(7{,}00\) п.п., а при 1200 К
остаются \(1{,}44\) п.п.
У статической модели авторский интервал разрешается даже при 1625 К.
Сохранение порядка до существенно более высоких температур в статическом
расчёте указывает на важность колебательной свободной энергии.

\begin{table}[htbp]
\centering\small
\caption{Верхние MC-сечения по авторскому C46 и независимой подгонке
при \(N=5120\). В каждой ячейке указаны \(x_\alpha/x_{\alpha_2}\),
ат.\% Al. Прочерк означает, что по данному сечению две границы не разрешены.}
\label{tab:ce-upper}
\begin{tabular}{@{}lrrr@{}}
\toprule
Модель & \(T\), К & Авторский C46 & Набор ATAT\\
\midrule
Статика & 1600 & \(22{,}89/24{,}16\) & \(23{,}13/24{,}82\)\\
        & 1625 & \(24{,}12/24{,}96\) & \(24{,}35/25{,}54\)\\
        & 1650 & --- & ---\\
HA      & 1200 & \(22{,}60/24{,}03\) & \(22{,}43/24{,}33\)\\
        & 1225 & --- & \(24{,}06/25{,}27\)\\
QHA     & 1200 & \(23{,}46/24{,}91\) & \(23{,}14/25{,}05\)\\
        & 1225 & --- & \(25{,}62/26{,}51\)\\
\bottomrule
\end{tabular}
\end{table}

Для семи уточнённых независимых сечений использовали по 69 окон с
шагом центров \(0{,}005\) по атомной доле Al. Разрешали внутренние
интервалы шире \(0{,}5\) п.п. и проверяли, что при переходе с
Ti-богатого конца к Al-богатому возрастает порядок D0$_{19}$.
В QHA при 1200 К замена наиболее медленного окна длинными
контрольными траекториями из двух начальных состояний дала
\(23{,}16/25{,}07\) ат.\% против \(23{,}14/25{,}05\) в таблице.
При 1225 К независимая QHA сохраняет интервал шириной \(0{,}89\) п.п.,
тогда как в авторских данных вблизи этой температуры две точки касания
не различаются. Порог ширины, шаг окон, конечный размер и зависимость
от подгонки ограничивают точность, с которой по этим сечениям можно
локализовать верхнее окончание двухфазной области.
